%
% !TEX root = Main.tex
\section{A simple robust parameter-free algorithm for stochastic \& adversarial rewards}
\label{s:UpperBOB}
%
In this section, we present a new learner 
and analyze its theoretical performance against any i.i.d.\, stochastic 
problem or any adversarially designed rewards.

We call the algorithm \ProbabilityONE{}, denote it by \Pone{}, %\todovout{is this english?, 'denote it'?},
and detail it in Figure~\ref{fig:probaoneAlgo}.
Intuitively,  \Pone{} pulls the estimated best arm with 
``probability'' one, the estimated second best arm with ``probability'' 
one half, the estimated third best arm with ``probability'' one 
third, and so on until pulling the estimated worst arm with ``probability'' one over $\nArms$.\textsuperscript{\ref{foot:break}} 
In order to get proper probabilities, we need to normalize  
them by the $\nArms$-th harmonic number %$\bar\log(\nArms)$ where %\todo{the algorithm could be defined so that the two best arm are pulled equally to make a bit more sense but I guess it doesn't really matter}
$\bar\log \nArms =\sum_{k=1}^\nArms (1/k)$, where $\bar\log \nArms\leq \log\nArms+1$ for all positive integers $\nArms$. \Pone{} uses a Zipf distribution over the estimated ranks
with an exponent of $1$~\citep{Powers98AE}.
% \todov{the relation between $\bar\log(\nArms)$ and $\log(\nArms)$ should be written somewhere}
% \todovout{michal, no? figure[H] for this algo? :) MV: like this?}
\begin{figure}[H]
	\centering
	\framebox{
		~~~~~~~~~    \begin{minipage}{.9\textwidth}
			%    Given: the number of arms $\nArms$ \\[.05cm]
			\vspace{.15cm}
			\textbf{For} $t=1,2, \ldots$\vspace{.1cm}
			
			$\qquad$\raisebox{.02cm}{\textcolor{bull}{$\blacktriangleright$}}~  Sort and rank the arms by decreasing 
			order of $\estCumulGainVector_{\cdot,t-1}$: Rank arm $k$ as $\tilde{\langle k \rangle}_t\in[K]$\footnote{Equalities between arms
				or comparisons with arms that have not been pulled 
				yet are broken arbitrarily.\label{foot:break}}.\vspace{.2cm}
			
			$\qquad$\raisebox{.02cm}{\textcolor{bull}{$\blacktriangleright$}}~  Select arm $\pulledArm_t\in [\nArms]$ with 
			$\learnerDist_{k,t}\triangleq\Pro\left(\pulledArm_t=k\mid\mathcal F_{t-1}\right) 
			\triangleq \displaystyle \frac{1}{~\tilde{\langle k \rangle}^{\vphantom{X}}_t~\bar\log\nArms~}$ 
			for all $k\in\setArms$. Here $\mathcal F_{t-1}$ is the information
			available before sampling, including the resolved ranking but excluding
			the fresh action draw and current rewards.\vspace{.2cm}
			
			\textbf{Recommend}, at any given round $t$,  $\recomArm_t \triangleq  \argmax_{k\in\setArms}~ 
			\estCumulGainVector_{k,t}$.
		\end{minipage}
	}
	\caption{The \ProbabilityONE{} (\Pone{}) algorithm}\label{fig:probaoneAlgo}
\end{figure}
\noindent
The estimate used in \Pone{} is 
$\estCumulGainVector_{k,t-1}$ for arm $k$ at round $t$. 
\Pone{} is heavily inspired  by Successive Rejects (\SR{}) of~\citet{Audibert10BA}, as both rank the
arms and attempt to allocate to arm~$k$ a number of pulls 
$\tcO\left( \timeHorizon/\langle k\rangle\right)$ according to its true rank $\langle k\rangle$. 
Our new learner is parameter-free and naturally anytime 
(agnostic of $\timeHorizon$).
% It only needs to know  $\nArms$. 
Like \SR{}, \Pone{} requires neither prior knowledge
nor estimation of the complexity parameters. 
Unlike \SR{}, \Pone{} does not divide the game into different 
sampling phases. \Pone{} follows the same deliberately simple rank-based
sampling rule at every round. %Therefore \Pone{} is a natural 
%anytime algorithm not needing to know the time horizon 
%$\timeHorizon$. 
%Additionally it makes the sampling sub-routine very 
%versatile and adaptive to other potential applications and settings. 



%\todom{if it is anytime, we remove $n$ from the algo\\ Victor: I removed it, but now  Is it confusing with the return $J_n$?}
As discussed in Section~\ref{s:FormuBOB}, in order to minimize the misidentification error in the stochastic case, it is crucial 
to limit 
the variance of the estimators for the best arms.
Therefore  \Pone{},  from
its very first pull, pulls more (i.e., with higher probability) the arms that are 
estimated to be among the best. 
First, this comes with almost no cost: Indeed, pulling 
the estimated best arm with probability 
$1/\bar\log\nArms$ does not prevent from pulling 
all the arms almost uniformly and more precisely with probability at 
least  $1/(\nArms\bar\log\nArms).$
Therefore, no suboptimal arm is actually left out in 
the early chase for the best arm, and the bounds $1/\learnerDist_{k,t}$
on the conditional variances are at most $\bar\log\nArms$ times
the corresponding bound under uniform sampling.
Second, it gives the learner the flexibility to adapt 
to different types of stochastic problems with 
different gap regimes.   If in a setup,  some arms 
are clearly suboptimal, it is helpful to pull the clearly best 
arm more and right from the beginning. This is a more flexible behavior than the one of algorithms that  are fixing the number of pulls for each arm during a fixed period  in advance. 

The key idea is that when \Pone{} ranks a promising arm
highly, it samples that arm more frequently,
reducing the variance of its importance-weighted
estimates. This, in turn, helps improve subsequent
rankings. We capture this feedback through virtual
comparison times that are introduced only for the
analysis and are not used by the algorithm.

Additionally, compared to a fixed-phase algorithm,  our
analysis is also more flexible:
We can analyze, for instance, the quality of the estimated ranking $\tilde{\langle\cdot\rangle}$ and therefore the adaptive sampling procedure of the arms 
at any round.
Actually, these rounds called \emph{comparison rounds}, 
can be chosen in a problem-dependent manner, in order to 
minimize the final probability of error. This is reflected in the complexity measure  present in the upper bound,
which is defined as the minimum over admissible choices
of comparison times.
\Pone{} remains agnostic to the optimal `virtual' phases,
which are introduced only for the analysis. We now define this 
new notion of complexity. First, we define  the proportion of rounds for 
comparison $\propTimeVec$  in a space~$\propTimeVecSpa.$
%
Let $\propTimeVecSpa\triangleq\left\{\propTimeVec
\in[0,1]^\nArms
: %\sum_{i=2}^\nArms \propTime_i \leq 1,
\timeHorizon\propTime_i\in\Integer ,  \forall i \in[\nArms],
1=\propTime_1=\propTime_2\geq\propTime_3\geq\ldots
\geq\propTime_{\nArms}>\propTime_{\nArms+1}=0
\right\}$. The complexity associated with the \Pone{} 
learner is $\complexityProOne$ and is defined 
first as
%\todov{explain what this is!}
%
%\todomi{the constants in the complexity term do not seem to match the one in the proof of Lemma 11}
\[
\complexityProOne(\propTimeVec)
\triangleq
\max_{k\in\setArms}
\frac{
	\sum_{i=\langle k \rangle}^{\nArms}
	(\propTime_{i}-\propTime_{i+1})
	i
	+
	\frac{1}{24}\nArms\propTime_{\langle k \rangle}\gap_{k}}
{\propTime^2_{\langle k \rangle}\gap^2_{k}}\bar\log\nArms
\quad\text{and then,}\quad
\complexityProOne
\triangleq
\min_{\propTimeVec\in\propTimeVecSpa}
\complexityProOne(\propTimeVec).
%
\]
%
%$\complexityProOne(\propTimeVec)$ has a similar shape to the variance complexity in~\cite{Gabillon11MB} (See a detailed discussion there). 
In  $\complexityProOne(\propTimeVec)$, for  arm $k$,  the term
$    \sum_{i=\langle k \rangle}^{\nArms}
(\propTime_{i}-\propTime_{i+1})
i$  corresponds to the sum of 
variances $\tcO(i)$, for $i\in[\langle k \rangle:\nArms]$,  during a proportion of time %\todom{this does not give a complete intuition of $\complexityProOne$}\todom{what is the variance of order $i$? may you mean by magnitude $i$? }
$\propTime_{i}-\propTime_{i+1}$ between the comparison 
rounds $i+1$ and~$i$. Indeed, 
%as discussed 
%later in the sketch of proof, 
on the good-ranking event described below, for $i<\nArms$ and any arm $k$ satisfying
$\mean_{(1)}-\mean_k<\gap_{(i+1)}/2$, we have
$\tilde{\langle k\rangle}_t\leq i$ during
$\timeHorizon\propTime_{i+1}<t\leq\timeHorizon\propTime_i$.
For $i=\nArms$, the same rank bound is automatic.
This leads, on these events and at the corresponding sampling rounds,
to $\learnerDist_{k,t}\geq1/(i\bar\log\nArms)$ and therefore, as
$\estGainVector_{k,t}=\gainVector_{k,t}\indicator{\pulledArm_t=k}/\learnerDist_{k,t}$,
to a conditional variance of $\estGainVector_{k,t}$ at most
$i\bar\log\nArms$.
In the denominator, $\propTime_{\langle k\rangle}$ specifies the
comparison time $\timeHorizon\propTime_{\langle k\rangle}$, the number
of elapsed rounds available before that comparison.
The additional term involving $\nArms/24$ comes from
the linear correction in the Bernstein--Freedman
concentration inequality.
%
\begin{restatable}[\textcolor{titleTh}{Upper bounds for \Pone{}}]{theorem}{lightres}    \label{th:UpBOB}
	For any stochastic  problem \stoPro{}
	with  complexity $\complexityProOne$
	and for any $\gainVector$ fixed by an oblivious adversary 
	with  complexity $\complexityUnif(\gainVector)$,
	%with $\gainVector_{k,t}\in [0,1]$ for all 
	%$t\in[\timeHorizon]$ and for all $k\in\setArms$,
	the  probabilities of error of~\Pone{}, denoted 
	$\ProErr_{\stoPro}(\timeHorizon)$ and $\ProErr_{\advPro(\gainVector)}(\timeHorizon)$
	in their respective environment, for any~$\timeHorizon$, simultaneously verify  	 
	%
	\[
	\ProErr_{\stoPro}(\timeHorizon)
	\leq
	2K^3n\exp\left(-\frac{\timeHorizon}{128\complexityProOne}\right)\quad \text{and}
	\quad
	\ProErr_{\advPro(\gainVector)}(\timeHorizon)
	\leq
	\nArms\exp\left(-\frac{3\timeHorizon}
	{40\bar\log(\nArms)\complexityUnif(\gainVector)}\right)\!\cdot
	\]
	%
\end{restatable}
%
%\todomi{on this page variance of the arms is used ambiguously or not clearly to describe the variance of the estimates }
\noindent\textbf{Sketch of the proof} 
\textit{(full proof in 
	Appendix~\ref{app:proofupBOB})}:
For the \textit{adversarial case}, it is enough that the learner pulls 
each arm with probability at least $1/(\nArms\bar\log\nArms)$ to obtain
the same complexity $\complexityUnif$ as \RULE,  up to a factor $\log\nArms$.
For the \textit{stochastic case},
we consider $\nArms-1$  arbitrary `virtual' 
phases that each  ends at round 
$\timeHorizon_i=\timeHorizon\propTime_i$. These comparison times are fixed
for the analysis and optimized only after deriving the final error bound.
Note that \Pone{} is oblivious to these values.
The phases are indexed in reverse order,
from $\nArms$ down to $2$.
Intuitively, $n_i$ is a round after which we expect the 
following event~$\eventbob_i$ to happen with  
high probability: 
%$\cO\left(\exp(-\timeHorizon/\complexityProOne)\right)$
For all integer $u$ with $\phaseTime_i<u\leq\timeHorizon+1$,
the ranking from estimates through $u-1$ satisfies
$\tilde{\langle k\rangle}_u\leq i-1<i$ whenever
$\mean_{(1)}-\mean_k<\gap_{(i)}/2$.
At $u=\timeHorizon+1$, this is the virtual final ranking consistent
with the recommendation.
The important consequence is that, on this event, any such arm $k$
is pulled with $\learnerDist_{k,t}\geq1/((i-1)\bar\log\nArms)$ at
sampling rounds $\phaseTime_i<t\leq\timeHorizon$, giving conditional
variance at most $(i-1)\bar\log\nArms$ in its estimates
$\estGainVector_{k,t}$.
Reducing these variances leads 
in turn to  better estimates in the 
rest of the game.
The proof works iteratively over the phases. We bound the probability
that $\eventbob_i$ fails together with the required earlier ranking
guarantees. On these good-ranking events, the conditional variance is
controlled, so martingale concentration bounds their joint occurrence
with a later comparison error. Summing these joint error probabilities
gives a bound on the probability of an incorrect final recommendation,
based on $\timeHorizon_2=\timeHorizon$ rounds of observations.

To bound $\ProErr_{\stoPro}(\timeHorizon)$,
we use the Bernstein inequality for martingale  differences. This inequality takes into account the 
variances and holds despite the dependencies of the random 
variables $\estGainVector_{k,t}$.

When minimizing $\complexityProOne$ by choosing $\propTime_k$, for $k\in\setArms$, we need to trade off between
short phases, possibly meaning not enough samples 
to discriminate the suboptimal arms (the denominator term of the suboptimal arms is small) and long phases,
which means that in early stages, the best arms are considered as badly 
ranked for a long time and the variance of their mean estimators is increasing 
with the length of the early phases (the numerator term for the best arms is larger).
%\todov{this previous paragraph should be written better}
%
\begin{restatable}{corollary}{cocores}\label{coco} %\todom{we have HBOB and HSTOBOB, aren't they the same? \\ Victor : Yes! Hbob is for stochastic setting \\ BUT WHY in THM 4 there is additionl superscript STO }
	For $\timeHorizon\geq\gap_{(\nArms)}/\gap_{(1)}$, the complexity
	$\complexityProOne$ of \Pone{} matches the complexity $\complexityBoth$
	from the lower bound of Theorem~\ref{th:lowBOB}
	up to log factors; more precisely,
	\[
	\complexityProOne\leq2\complexityBoth(\bar\log\nArms)^2
	=\cO\left(\complexityBoth\log^2\nArms\right)\!.
	\]
	For every positive integer horizon $\timeHorizon$, the stochastic
	probability of error of \Pone{} satisfies
	\[
	\ProErr_{\stoPro}(\timeHorizon)
	\leq2\nArms^3\timeHorizon\exp\left(
	-\frac{\timeHorizon}{256\complexityBoth(\bar\log\nArms)^2}
	\right)\!.
	\]
\end{restatable}
%\todom{this setting of a minimizes P1}
\noindent The result of Corollary~\ref{coco} is obtained by setting
$\propTime_k=\left\lceil\timeHorizon\gap_{(1)}/\gap_{(k)}\right\rceil/\timeHorizon,$
$\forall k\in\setArms$.
We round each ideal comparison time upward to the next round.
This changes only the analysis, leaving the learner and its horizon
unchanged, and costs at most a factor two in the complexity in the
stated budget range. For smaller budgets, the displayed probability
bound is trivial. Notice that the same gap ratio also guides the lower-bound
construction of Theorem~\ref{th:lowBOB}. For its reference gaps,
we use $\dividefac_i=\min\{1,2\gap_1/\gap_i\}$ to leave enough
margin in the adversarial gaps after allowing for random
fluctuations. The full proof  of Corollary~\ref{coco}
is in Appendix~\ref{s:moreOnRem}, 
where we also discuss the relation between 
$\complexityProOne$ and $\complexitySR$ and 
$\complexityBoth$ for different regimes of the 
gaps.
Corollary~\ref{coco} demonstrates that \Pone{} achieves the best that can be wished 
for in the two worlds, up to log factors.
%
%
%\todov{Michal, do you agree with the following remark 8?}
% \todomi{Let's discuss the detail.
% First, what happens to the STO part? 
%Second, why in the ADV 1/2 of P1 and 1/2 of RULE, gives you this?}
\begin{remark}
	In the adversarial case, a  modification to $\Pone{}$ could lead to a similar upper bound as for \RULE{},  where  $\complexityUnif(\gainVector)$ appears instead of $\complexityUnif(\gainVector)\bar\log \nArms$.
	At each round, with probability one half we can play according to \RULE{} and otherwise use \Pone{}. The importance-weighted estimator uses the actual sampling probability $\learnerDist_{k,t}$ of this mixture. We keep the  recommendation $\recomArm_\timeHorizon=\argmax_{k\in\setArms}~ 
	\estCumulGainVector_{k,\timeHorizon}$.
	Such a mixture could retain useful stochastic adaptation.
\end{remark}
%
%\todov{we will probably have to remove the following remark}
\begin{remark}
	We studied the hard (adversarial) and easy data (stochastic) settings. 
	However, as discussed by~\cite{Seldin14OP}~and~\cite{Allesiardo17NS}, we can consider intermediate settings of difficulty.
	First, one could investigate whether the result in the stochastic case holds up
	to constants when the gaps of the arms do not change by more than the 
	same numerical factor during the game.
	More generally, we could design variants of $\complexityProOne$ under 
	the assumption that after some round $t'$ the (ground-truth) ranks of all the 
	arms are upper bounded each by a constant. %\todom{I changed to upper bounded each by a constant. OK victor ?}
	A variant of \Pone{} might adapt after~$t'$ and rank every arm at most according to its 
	maximum rank. Such results might even hold in 
	a case of a change of the identity of the best arm in the game.
	%
	%Then a careful reader will se that one can easily  redefine the notion of gaps and have a notion of maximal ranking to extend our results. 
	%We look at a version of the problem where the best arm is define as in the stochastic problem, just not necessarily stochastic.
	%So we have a distribution at each time and arm $\nu_{k,t}$ and a best arm at each time step $\bestArm_{t}=\argmax_{k\in\setArms}\Exp\left[\sum_{t'=1}^t \mean_{k}(t)\right]$
	%\todov{is this well defined?}
	%If there exist a time $t_s$, maximal rankings $r_{k,s}$ and minimal gaps such that for all times $t$, $t\geq t_s=\propTime_s\timeHorizon$, $\propTime_s\in[0,1]$ we have that for all $k\in\setArms$ both $\langle k \rangle_t \leq r_{k,s}$ and
	%$\langle k \rangle_t \leq r_{k,s}$  
	%
\end{remark}%